\section{Relevancia del Estudio}

\subsection{Impacto potencial del estudio en la teoría matemática}
El impacto principal de esta investigación radica en su enriquecimiento teórico al campo de la Optimización Continua y el Análisis Convexo aplicado a grafos. Al trascender el paradigma de las aproximaciones lineales, el estudio proporciona demostraciones analíticas que clarifican el comportamiento geométrico de regiones factibles limitadas por inecuaciones no lineales. Específicamente, el escrutinio sobre la matriz Hessiana de las funciones de costo y el análisis de sensibilidad de los multiplicadores de Lagrange (en el marco de las condiciones de Karush-Kuhn-Tucker) ofrecerá nuevas perspectivas matemáticas sobre cómo las no linealidades originadas por economías de escala afectan la compacidad del espacio de soluciones y la unicidad de los óptimos globales.

\subsection{Contribución esperada del estudio a la resolución de problemas matemáticos más amplios}
El aparato deductivo y los resultados algorítmicos desarrollados en esta tesis trascenderán la particularidad del grafo logístico y servirán de cimiento para el abordaje de problemas matemáticos de escala macro. La evaluación teórica de la complejidad temporal y la cota del error numérico en algoritmos de segundo orden frente a matrices dispersas (como la matriz de incidencia del grafo) arroja luz sobre la eficiencia computacional de los Métodos de Puntos Interiores en sistemas multivariables. Estas demostraciones son altamente valoradas y extrapolables para la resolución de la clase abstracta de problemas de "Flujo de Múltiples Mercancías con Costo No Lineal" (\textit{Non-linear Multi-commodity Flow Problems}), cuyo estudio de intratabilidad (espectro NP-Hard) es un área hiperactiva en las matemáticas computacionales.

\subsection{Posibles aplicaciones de los resultados en otras áreas de la matemática o en campos relacionados}
Debido a que el modelo matemático se construye axiomáticamente sobre una estructura topológica abstracta $G=(V,E)$, las ecuaciones, teoremas y el algoritmo iterativo subyacente pueden ser despojados de su contexto logístico para solucionar problemas isomórficos en disciplinas afines:
\begin{itemize}
    \item \textbf{Teoría de Redes de Telecomunicaciones:} El ruteo óptimo de paquetes de datos a través de servidores de internet, donde la congestión del tráfico genera retardos exponenciales (costo no lineal), se rige por las mismas ecuaciones de conservación de flujo y capacidad diferencial modeladas en esta tesis.
    \item \textbf{Mecánica de Fluidos e Hidráulica de Redes:} El cálculo de caudales óptimos en sistemas de gasoductos o acueductos interregionales. En estos sistemas, el principio de continuidad (nodos) y las ecuaciones de pérdida de carga por fricción (arcos) constituyen intrínsecamente un problema no lineal equivalente.
    \item \textbf{Ciencia de Datos y Machine Learning:} El estudio algorítmico del descenso por gradientes en estructuras de red es el soporte matemático fundamental de las Redes Neuronales de Grafos (\textit{Graph Neural Networks - GNN}), permitiendo que las estrategias de minimización desarrolladas aquí sean adaptables a arquitecturas de aprendizaje profundo.
\end{itemize}
